\section*{Chapter \ref{ch:b2:curves} --- Curves}

\begin{solution}{exo:b2:curves:1}
$\gamma'(t) = (1 - \cos t,\ \sin t)$, so
\[
\norm{\gamma'(t)}^2 = (1 - \cos t)^2 + \sin^2 t
= 2 - 2\cos t = 4\sin^2\tfrac t2 ,
\]
and $\norm{\gamma'(t)} = 2\sin\frac t2$ (nonnegative on $[0,
2\pi]$). Hence
\[
L = \int_0^{2\pi} 2\sin\tfrac t2\,\dd t
= \Bigl[-4\cos\tfrac t2\Bigr]_0^{2\pi} = 8 :
\]
one arch of the cycloid has \omterm{def:b2:curves:length}{length} $8$ (for a wheel of radius
$1$) --- a famous result of Wren, with no $\pi$ in sight.
\end{solution}

\begin{solution}{exo:b2:curves:2}
With $x = a\cos t$, $y = b\sin t$: $x' = -a\sin t$, $y' = b\cos
t$, $x'' = -a\cos t$, $y'' = -b\sin t$, so by
\cref{prop:b2:curves:kappaformula}
\[
\kappa(t)
= \frac{x'y'' - y'x''}{(x'^2 + y'^2)^{3/2}}
= \frac{ab\sin^2 t + ab\cos^2 t}
       {(a^2\sin^2 t + b^2\cos^2 t)^{3/2}}
= \frac{ab}{(a^2\sin^2 t + b^2\cos^2 t)^{3/2}} .
\]
The denominator is minimal when $\sin t = 0$ (value $b^3$, points
$(\pm a, 0)$) and maximal when $\cos t = 0$ (value $a^3$, points
$(0, \pm b)$), since $a > b$. Hence $\kappa$ is maximal at the
ends of the major axis, $\kappa_{\max} = a/b^2$, and minimal at
the ends of the minor axis, $\kappa_{\min} = b/a^2$: the ellipse
bends most sharply at the tips of its long axis.
\end{solution}

\begin{solution}{exo:b2:curves:3}
For the graph $\gamma(x) = (x, \cosh x)$: $\norm{\gamma'(x)} =
\sqrt{1 + \sinh^2 x} = \cosh x$, so the \omterm{def:b2:curves:length}{arc length} from $0$ to $x$
is $\int_0^x \cosh u\,\dd u = \sinh x$. \omterm{thm:b2:curves:frenet2d}{Curvature} of a graph
(\cref{prop:b2:curves:kappaformula}):
\[
\kappa(x) = \frac{f''(x)}{(1 + f'(x)^2)^{3/2}}
= \frac{\cosh x}{\cosh^3 x} = \frac{1}{\cosh^2 x} ,
\]
so $R(x) = \cosh^2 x$, as announced. Note the neat coincidence
$R(x) = 1 + s(x)^2$ with $s = \sinh x$ the \omterm{def:b2:curves:length}{arc length}: the
catenary's \omterm{def:b2:curves:curvature}{radius of curvature} grows with the square of the \omterm{def:b2:curves:length}{arc
length} from the vertex.
\end{solution}

\begin{solution}{exo:b2:curves:4}
$\gamma'(t) = e^t(\cos t - \sin t,\ \sin t + \cos t)$, so
\[
\langle \gamma(t), \gamma'(t)\rangle = e^{2t}
\bigl(\cos t(\cos t - \sin t) + \sin t(\sin t + \cos t)\bigr)
= e^{2t},
\]
while $\norm{\gamma(t)} = e^t$ and $\norm{\gamma'(t)} =
e^t\sqrt 2$. Hence
\[
\cos\angle\bigl(\gamma, \gamma'\bigr)
= \frac{e^{2t}}{e^t \cdot e^t\sqrt2} = \frac{1}{\sqrt2} :
\]
the tangent always makes the angle $\pi/4$ with the radius --- the
\emph{equiangular} property of the logarithmic spiral. \omterm{def:b2:curves:length}{Arc length}
on $(-\infty, 0]$:
\[
\int_{-\infty}^0 \norm{\gamma'(t)}\,\dd t
= \sqrt2\int_{-\infty}^0 e^t\,\dd t = \sqrt 2 ,
\]
finite although the spiral winds infinitely many times around the
origin. \omterm{thm:b2:curves:frenet2d}{Curvature}: with $x'y'' - y'x''$ computed from $\gamma'' =
e^t(-2\sin t,\ 2\cos t)$,
\[
x'y'' - y'x'' = e^{2t}\bigl(2\cos t(\cos t - \sin t)
+ 2\sin t(\sin t + \cos t)\bigr) = 2e^{2t},
\]
so $\kappa(t) = \dfrac{2e^{2t}}{(e^t\sqrt2)^3} =
\dfrac{1}{e^t\sqrt2}$: the \omterm{thm:b2:curves:frenet2d}{curvature} is $1/(\sqrt2\,
\norm{\gamma})$, decaying as the spiral grows.
\end{solution}

\begin{solution}{exo:b2:curves:5}
\emph{First arc:} $\gamma(t) = (t^2,\ t^4 + t^5)$. Derivatives at
$0$: $\gamma'' = (2, 0) \neq 0$, so $p = 2$. Then $\gamma^{(3)}(0)
= (0, 0)$, $\gamma^{(4)}(0) = (0, 24)$, not collinear with $(2,
0)$: $q = 4$. Both even: \emph{cusp of the second kind} --- both
branches leave in the direction $+u = (1,0)$ and stay on the same
side of the tangent. (Indeed $y = x^2 \pm x^{5/2}$ on the two
branches: same sign for small $x$.)

\emph{Second arc:} $\gamma(t) = (t^3, t^4)$. $\gamma'(0) =
\gamma''(0) = 0$, $\gamma^{(3)}(0) = (6, 0)$: $p = 3$, odd. Next
$\gamma^{(4)}(0) = (0, 24)$: $q = 4$, even. Odd--even:
\emph{ordinary point} --- despite the vanishing velocity, the
trajectory $y = x^{4/3}$ crosses the origin smoothly, staying
above its tangent $y = 0$.
\end{solution}

\begin{solution}{exo:b2:curves:6}
Differentiate $c(s) = \gamma(s) + \dfrac{1}{\kappa(s)}N(s)$ using
the plane Frenet formulas
(\cref{thm:b2:curves:frenet2d}):
\[
c'(s) = T(s) - \frac{\kappa'(s)}{\kappa(s)^2}N(s)
+ \frac{1}{\kappa(s)}\,\bigl(-\kappa(s)T(s)\bigr)
= -\frac{\kappa'(s)}{\kappa(s)^2}\,N(s) ,
\]
the tangent terms cancelling exactly. So the velocity of the
evolute is carried by $N(s)$, which directs the normal line of
$\gamma$ at $\gamma(s)$ --- and the point $c(s)$ lies on that very
normal line: the evolute is the \emph{\omterm{pb:b2:curves:1}{envelope} of the normals}.
(Where $\kappa' = 0$ the evolute has a singular point; this is
what produces the cusps of the evolute of an ellipse.)
\end{solution}

\begin{solution}{exo:b2:curves:7}
Let $G(s) = F(s)F(s)^{\mathsf T}$. Then, using $F' = FA$ and
$(F^{\mathsf T})' = (F')^{\mathsf T} = A^{\mathsf T}F^{\mathsf T}$,
\[
G' = F'F^{\mathsf T} + F(F^{\mathsf T})'
= FAF^{\mathsf T} + FA^{\mathsf T}F^{\mathsf T}
= F(A + A^{\mathsf T})F^{\mathsf T} = 0
\]
by antisymmetry. So $G$ is constant on the interval, equal to
$G(s_0) = F(s_0)F(s_0)^{\mathsf T} = I$: $F(s)$ is orthogonal for
every $s$. (Alternatively, without computing $G'$ to zero: both
$G$ and the constant $I$ solve the linear system $Y' = YA +
A^{\mathsf T}Y$ with the same initial value, and Cauchy--Lipschitz
uniqueness for linear systems, \cref{ch:b2:diffeq}, forces $G
\equiv I$.)

\emph{Relevance:} the Frenet system $(T, N, B)' = (T, N, B)\,A(s)$
has the antisymmetric coefficient matrix
\[
A = \begin{pmatrix} 0 & -\kappa & 0\\ \kappa & 0 & -\tau\\
0 & \tau & 0\end{pmatrix}
\]
(columns expressing $T', N', B'$). The computation above shows
that a solution frame that starts orthonormal \emph{stays}
orthonormal --- the key step in the fundamental theorem
reconstructing a curve from $(\kappa, \tau)$.
\end{solution}

\begin{solution}{exo:b2:curves:8}
By \cref{thm:b2:curves:fundamental} (uniqueness part), there is a
$\mathcal{C}^1$ angle function $\varphi$ with $T(s) =
(\cos\varphi(s), \sin\varphi(s))$ and $\varphi' = \kappa$. Hence
\[
\int_0^L \kappa(s)\,\dd s = \varphi(L) - \varphi(0) .
\]
Since the arc is closed of period $L$, $T(L) = T(0)$:
$(\cos\varphi(L), \sin\varphi(L)) = (\cos\varphi(0),
\sin\varphi(0))$, so $\varphi(L) - \varphi(0) \in 2\pi\Z$. The
total \omterm{thm:b2:curves:frenet2d}{curvature} of a closed curve is therefore always an integer
multiple of $2\pi$ --- the integer being the \emph{winding number}
of the tangent (the number of full turns $T$ makes). For the
circle of radius $R$: $\kappa = 1/R$ and $L = 2\pi R$, total
\omterm{thm:b2:curves:frenet2d}{curvature} $2\pi$, winding number $1$; the theorem of turning
tangents states this value holds for every simple closed curve.
\end{solution}

\begin{solution}{exo:b2:curves:9}
Since $\tau \equiv 0$, the curve lies in a plane
(\cref{prop:b2:curves:torsion}); work in that plane. Consider the
candidate center
\[
c(s) = \gamma(s) + \frac{1}{\kappa}N(s)
\qquad (\kappa \text{ constant}).
\]
Differentiating with the Frenet formulas ($N' = -\kappa T + \tau B
= -\kappa T$ here):
\[
c'(s) = T + \frac1\kappa(-\kappa T) = 0 ,
\]
so $c$ is a constant point $\Omega$. Then $\norm{\gamma(s) -
\Omega} = \norm{-\frac1\kappa N(s)} = \frac1\kappa$ for all $s$:
the curve lies on the circle of center $\Omega$ and radius
$1/\kappa$ (in its plane), and being a nonconstant arc of it, it
is an arc of that circle.
\end{solution}

\begin{solution}{exo:b2:curves:10}
For the graph $\gamma(x) = (x, x^2/2)$, $\norm{\gamma'(x)} =
\sqrt{1 + x^2}$, so $L = \int_0^a\sqrt{1+x^2}\,\dd x$.
Substituting $x = \sinh u$ ($\dd x = \cosh u\,\dd u$, $u$ from
$0$ to $u_a = \ln(a + \sqrt{1+a^2})$):
\[
L = \int_0^{u_a}\cosh^2 u\,\dd u
= \frac12\bigl[u + \sinh u\cosh u\bigr]_0^{u_a}
= \frac12\Bigl(\ln\bigl(a + \sqrt{1+a^2}\bigr) +
a\sqrt{1+a^2}\Bigr),
\]
using $\cosh^2 u = \frac{1 + \cosh 2u}2$ and $\sinh u_a = a$,
$\cosh u_a = \sqrt{1 + a^2}$.
\end{solution}

\begin{solution}{exo:b2:curves:11}
Parametrize by \omterm{def:b2:curves:length}{arc length}
(\cref{thm:b2:curves:arclength}) and set $\lambda(s) =
\langle P - \gamma(s), T(s)\rangle$, a $\mathcal C^1$
function; since $P$ lies on the \omterm{def:b2:curves:arc}{tangent line} at $\gamma(s)$,
the vector $P - \gamma(s)$ is collinear with $T(s)$, so $P =
\gamma(s) + \lambda(s)T(s)$. Differentiating,
\[
0 = T(s) + \lambda'(s)T(s) + \lambda(s)T'(s)
= \bigl(1 + \lambda'(s)\bigr)T(s) + \lambda(s)T'(s),
\]
and $T'(s) \perp T(s)$ (differentiate $\norm T^2 = 1$), so
both components vanish: $\lambda' = -1$ and $\lambda T' = 0$.
Then $\lambda(s) = c - s$ vanishes at most once, so $T' = 0$
on a dense set, hence everywhere by \omterm{def:b2:metric:continuity}{continuity}: $T$ is a
constant unit vector and $\gamma(s) = \gamma(s_0) + (s -
s_0)T$: a straight line (through $P$, as it must be).
\end{solution}

\begin{solution}{exo:b2:curves:12}
(a) The Frenet matrix
\[
A(s) = \begin{pmatrix} 0 & -\kappa & 0\\ \kappa & 0 & -\tau\\
0 & \tau & 0\end{pmatrix}
\]
has \omterm{def:b2:metric:continuity}{continuous} entries, so the linear system $F' = FA(s)$,
$F(s_0) = F_0$ (a direct orthonormal matrix) has a unique
solution on all of $J$ (\cref{thm:b2:diffeq:linear}). By
\cref{exo:b2:curves:7}, $F(s)$ is orthogonal for every $s$;
$\det F$ is \omterm{def:b2:metric:continuity}{continuous} with values in $\{\pm1\}$ and equals
$1$ at $s_0$, so $F(s)$ is direct for all $s$.

(b) Read off the rows $T, N, B$ of $F$ (so that $T' = \kappa
N$, $N' = -\kappa T + \tau B$, $B' = -\tau N$) and set
$\gamma(s) = \gamma_0 + \int_{s_0}^s T(u)\,\dd u$. Then
$\gamma' = T$ is a unit vector: unit speed; $T' = \kappa N$
with $\kappa > 0$ and $N$ unit orthogonal to $T$, so $\gamma$
is biregular with \omterm{thm:b2:curves:frenet2d}{curvature} $\norm{T'} = \kappa$ and
principal normal $N$; the binormal is $T \wedge N = B$
(direct orthonormal frame), and $B' = -\tau N$ identifies the
\omterm{thm:b2:curves:frenet3d}{torsion} as $\tau$.

(c) Let $\gamma_1, \gamma_2$ be unit-speed biregular curves
with the same $(\kappa, \tau)$. There is a unique direct
isometry $\Phi = \rho + w$ ($\rho \in SO(3)$) sending
$\gamma_1(s_0)$ to $\gamma_2(s_0)$ and the Frenet frame of
$\gamma_1$ at $s_0$ to that of $\gamma_2$ at $s_0$. The curve
$\Phi\circ\gamma_1$ is unit-speed with the same invariants
(its frame is $\rho$ applied to that of $\gamma_1$, and
$\rho$ preserves cross products, being direct). Now the
frames of $\Phi\circ\gamma_1$ and $\gamma_2$ both solve $F' =
FA(s)$ with the same initial value, so they coincide by
uniqueness; in particular the tangents agree, and integrating
from the common point $s_0$: $\Phi\circ\gamma_1 = \gamma_2$.
\end{solution}

\begin{solution}{pb:b2:curves:1}
\textbf{1.} The system is linear in $(x, y)$ with \omterm{def:b2:linalg:det}{determinant}
$\Delta(t) = a b' - a'b \neq 0$: Cramer's rule gives the
unique solution
\[
x = \frac{c b' - c' b}{\Delta}, \qquad
y = \frac{a c' - a' c}{\Delta}.
\]

\textbf{2.} Since $a(t)x(t) + b(t)y(t) = c(t)$ identically,
differentiating gives $a'x + b'y + ax' + by' = c'$; the
second characteristic equation kills $a'x + b'y - c'$, so
$a(t)x'(t) + b(t)y'(t) = 0$: $E'(t)$ is orthogonal to $(a,
b)$, hence parallel to $(-b, a)$, the direction of $D_t$. As
$E(t) \in D_t$ (first equation) and $E'(t) \neq 0$, the line
$D_t$ is exactly the \omterm{def:b2:curves:arc}{tangent line} of the curve $E$ at $E(t)$.

\textbf{3.} Here $(a, b, c) = (t, -1, t^2/2)$, so $\Delta =
t\cdot0 - 1\cdot(-1) = 1$ and
\[
x = \frac{c b' - c' b}{\Delta} = \tfrac{t^2}2\cdot 0 +
t = t, \qquad
y = \frac{a c' - a' c}{\Delta} = t\cdot t - \tfrac{t^2}2 =
\tfrac{t^2}2 :
\]
the \omterm{pb:b2:curves:1}{envelope} of the \omterm{def:b2:curves:arc}{tangent lines} of the parabola is the
parabola, as it should be.

\textbf{4.} The normal line is $\langle M, T(s)\rangle =
\langle\gamma(s), T(s)\rangle$: coefficients $a = T_1$, $b =
T_2$, $c = \langle\gamma, T\rangle$. Differentiating with
Frenet ($T' = \kappa N$): $a' = \kappa N_1$, $b' = \kappa
N_2$, and $c' = \langle T, T\rangle + \langle\gamma, \kappa
N\rangle = 1 + \kappa\langle\gamma, N\rangle$. The second
characteristic equation $\kappa\langle M, N\rangle = 1 +
\kappa\langle\gamma, N\rangle$ reads $\kappa\langle M -
\gamma, N\rangle = 1$. The first says $M - \gamma \perp T$,
so $M - \gamma = \mu N$ with $\mu = 1/\kappa$: the
characteristic point is $\gamma + \frac1\kappa N$, the \omterm{def:b2:curves:curvature}{center
of curvature}, and the \omterm{pb:b2:curves:1}{envelope} of the normals is the evolute
of \cref{exo:b2:curves:6}. Finally $\Delta = T_1\kappa N_2 -
\kappa N_1 T_2 = \kappa\det(T, N) = \kappa \neq 0$.

\textbf{5.} Pencil: the system $x\cos\theta + y\sin\theta =
0$, $-x\sin\theta + y\cos\theta = 0$ has \omterm{def:b2:linalg:det}{determinant} $1$ and
solution $(0,0)$ for every $\theta$: $E \equiv (0,0)$, $E'
\equiv 0$, and there is no curve --- merely the common point
of all the lines. Parallel family: $a' = b' = 0$ gives
$\Delta \equiv 0$ and the second equation $0 = c'(t)$, which
fails as soon as the family actually moves: no characteristic
point, and indeed a family of parallel lines touches no
curve along all its members.

\textbf{6.} The line through $(\cos t, 0)$ and $(0, \sin t)$
is $\frac x{\cos t} + \frac y{\sin t} = 1$, i.e.\ $x\sin t +
y\cos t = \sin t\cos t$. With $(a, b, c) = (\sin t, \cos t,
\sin t\cos t)$: $a' = \cos t$, $b' = -\sin t$, $c' = \cos
2t$, $\Delta = -\sin^2 t - \cos^2 t = -1$. Cramer:
\begin{align*}
x &= \frac{cb' - c'b}{-1} = \sin^2 t\cos t + \cos 2t\cos t
= \cos t\,(\sin^2 t + \cos^2 t - \sin^2 t) = \cos^3 t,\\
y &= \frac{ac' - a'c}{-1} = \sin t\cos^2 t - \sin t\cos 2t
= \sin t\,(\cos^2 t - \cos^2 t + \sin^2 t) = \sin^3 t .
\end{align*}
And $(\cos^3t)^{2/3} + (\sin^3t)^{2/3} = 1$: the \omterm{pb:b2:curves:1}{astroid}.

\textbf{7.} $E'(t) = 3(-\cos^2 t\sin t,\ \sin^2 t\cos t)$
vanishes at $t = 0$. There, $E''(0) = (-3, 0) \neq 0$ gives
$p = 2$; the $x$-component of $E$ is even in $t$, so
$E'''(0) = (0, 6)$, not collinear: $q = 3$. Even--odd: cusp
of the first kind at $(1, 0)$
(\cref{prop:b2:curves:local}), with tangent along the
$x$-axis. The symmetries $x \mapsto -x$, $y \mapsto -y$, $(x,
y) \mapsto (y, x)$ of the \omterm{pb:b2:curves:1}{astroid} transport the cusp to $(-1,
0)$ and $(0, \pm1)$.

\textbf{8.} $E'(t) = 3\sin t\cos t\,(-\cos t, \sin t)$, so
$\norm{E'(t)} = 3\abs{\sin t\cos t} = \tfrac32\abs{\sin 2t}$.
One quadrant: $\int_0^{\pi/2}\tfrac32\sin 2t\,\dd t =
\tfrac32$, and by symmetry the total \omterm{def:b2:curves:length}{length} is $4 \cdot
\tfrac32 = 6$.

\textbf{9.} $E(t) - P_t = (\cos^3 t - \cos t,\ \sin^3 t) =
\sin^2 t\,(-\cos t,\ \sin t) = \sin^2 t\,(Q_t - P_t)$. So the
contact point is the \omterm{def:b2:affine:barycenter}{barycenter} of $(P_t, \cos^2 t)$ and
$(Q_t, \sin^2 t)$: as $t$ runs from $0$ to $\pi/2$ it slides
from the floor end to the wall end of the ladder.

\textbf{10.} In the first quadrant the region under the
\omterm{pb:b2:curves:1}{astroid} has area $\int_0^1 y\,\dd x$ with $x = \cos^3 t$
decreasing from $1$ to $0$ as $t$ goes from $0$ to $\pi/2$:
\[
\int_0^1 y\,\dd x
= \int_{\pi/2}^{0}\sin^3 t\,(-3\cos^2 t\sin t)\,\dd t
= 3\int_0^{\pi/2}\sin^4 t\cos^2 t\,\dd t .
\]
Linearize: $\sin^4 t\cos^2 t = (\sin t\cos t)^2\sin^2 t =
\tfrac18\bigl(\sin^2 2t - \sin^2 2t\cos 2t\bigr)$, and
$\int_0^{\pi/2}\sin^2 2t\,\dd t = \tfrac\pi4$ while
$\int_0^{\pi/2}\sin^2 2t\cos 2t\,\dd t =
\bigl[\tfrac{\sin^3 2t}6\bigr]_0^{\pi/2} = 0$. So the
integral is $\tfrac\pi{32}$, the quadrant area
$\tfrac{3\pi}{32}$, and the enclosed area $4 \cdot
\tfrac{3\pi}{32} = \tfrac{3\pi}8$.

\textbf{11.} The tangent at $\gamma(t) = (t, t^2/2)$ is
directed by $(1, t)$, so the normal line is $\{(x, y) : (x -
t) + t\,(y - t^2/2) = 0\}$, i.e.\ $x + t\,y = t +
\tfrac{t^3}2$.

\textbf{12.} $(a, b, c) = (1, t, t + t^3/2)$: $a' = 0$, $b' =
1$, $c' = 1 + \tfrac32 t^2$, $\Delta = 1$. Then $y = ac' -
a'c = 1 + \tfrac32t^2$ and
\[
x = cb' - c'b = t + \tfrac{t^3}2 - t\Bigl(1 +
\tfrac32t^2\Bigr) = -t^3 .
\]
Eliminating $t$: $t^2 = \tfrac23(y - 1)$ and $x^2 = t^6 =
\tfrac8{27}(y - 1)^3$: a semicubical parabola with vertex
$(0, 1)$.

\textbf{13.} $T = (1, t)/\sqrt{1+t^2}$, $N = (-t,
1)/\sqrt{1+t^2}$, and $\kappa = (1+t^2)^{-3/2}$, so
\[
\gamma + \frac1\kappa N = (t, \tfrac{t^2}2) +
(1+t^2)\,(-t, 1) = \Bigl(-t^3,\ 1 + \tfrac32t^2\Bigr),
\]
the same curve: the \omterm{pb:b2:curves:1}{envelope} of the normals is the locus of
the centers of \omterm{thm:b2:curves:frenet2d}{curvature}, as question 4 promised.

\textbf{14.} $E'(t) = (-3t^2, 3t)$ vanishes at $t = 0$;
$E''(0) = (0, 3) \neq 0$ gives $p = 2$ and $E'''(0) = (-6,
0)$ gives $q = 3$: a cusp of the first kind at $(0, 1)$. The
vertex is where $\kappa = (1+t^2)^{-3/2}$ is maximal, so
$\kappa'(0) = 0$ and the evolute's velocity
$-\frac{\kappa'}{\kappa^2}N$ vanishes exactly there: cusps of
the evolute sit at the extrema of \omterm{thm:b2:curves:frenet2d}{curvature}.

\textbf{15.} The normal at parameter $t$ passes through
$(x_0, y_0)$ iff $x_0 + t\,y_0 = t + \tfrac{t^3}2$, i.e.
\[
\frac{t^3}2 + (1 - y_0)\,t - x_0 = 0 ,
\]
a cubic in $t$: one or three real roots (counted without
multiplicity, for generic points). On the axis $x_0 = 0$ it
factors as $t\bigl(\tfrac{t^2}2 + 1 - y_0\bigr) = 0$: the
root $t = 0$ (the axis is the normal at the vertex), plus
$t = \pm\sqrt{2(y_0 - 1)}$ when $y_0 > 1$. So: one normal
for $y_0 < 1$, three for $y_0 > 1$, and at $y_0 = 1$ the
triple root marks the cusp of the evolute. In general a
double root of the cubic means the point satisfies both the
line equation and its $t$-derivative --- it lies \emph{on the
\omterm{pb:b2:curves:1}{envelope}}: the evolute is precisely the boundary between the
one-normal and three-normal regions.

\textbf{16.} Reflection in the mirror reverses the normal
component of the direction and keeps the tangential one:
writing $u = \langle u, n\rangle n + u_{\mathrm{tan}}$, the
reflected direction is $u_{\mathrm{tan}} - \langle u,
n\rangle n = u - 2\langle u, n\rangle n$. Here $\langle u,
n\rangle = \cos\theta$, so
\[
v = (1, 0) - 2\cos\theta\,(\cos\theta, \sin\theta)
= (1 - 2\cos^2\theta,\ -2\sin\theta\cos\theta)
= -(\cos2\theta,\ \sin2\theta).
\]

\textbf{17.} The reflected ray passes through $P_\theta =
(\cos\theta, \sin\theta)$ with direction $(\cos2\theta,
\sin2\theta)$; a normal vector is $(-\sin2\theta,
\cos2\theta)$, so the line is
\[
-\sin2\theta\,(x - \cos\theta) + \cos2\theta\,(y -
\sin\theta) = 0,
\]
and the constant is $-\sin2\theta\cos\theta +
\cos2\theta\sin\theta = -\sin\theta$: multiplying by $-1$,
$x\sin2\theta - y\cos2\theta = \sin\theta$.

\textbf{18.} $(a, b, c) = (\sin2\theta, -\cos2\theta,
\sin\theta)$: $a' = 2\cos2\theta$, $b' = 2\sin2\theta$, $c' =
\cos\theta$, $\Delta = 2\sin^22\theta + 2\cos^22\theta = 2$.
Cramer, then product-to-sum formulas:
\begin{align*}
x &= \frac{2\sin\theta\sin2\theta +
\cos\theta\cos2\theta}{2}
= \frac{(\cos\theta - \cos3\theta) + \frac12(\cos\theta +
\cos3\theta)}{2} = \frac{3\cos\theta - \cos3\theta}4,\\
y &= \frac{\sin2\theta\cos\theta - 2\cos2\theta\sin\theta}2
= \frac{\frac12(\sin3\theta + \sin\theta) - (\sin3\theta -
\sin\theta)}2 = \frac{3\sin\theta - \sin3\theta}4 :
\end{align*}
the \omterm{pb:b2:curves:1}{nephroid}, a closed curve with two cusps.

\textbf{19.} Differentiating and factoring with $\sin3\theta
- \sin\theta = 2\cos2\theta\sin\theta$, $\cos\theta -
\cos3\theta = 2\sin2\theta\sin\theta$:
\[
E'(\theta) = \tfrac34\bigl(\sin3\theta - \sin\theta,\
\cos\theta - \cos3\theta\bigr)
= \tfrac32\sin\theta\,(\cos2\theta, \sin2\theta),
\]
parallel to the reflected direction of question 16: each
reflected ray is tangent to the caustic, as the \omterm{pb:b2:curves:1}{envelope}
property demands. $E' = 0$ exactly at $\sin\theta = 0$:
$E(0) = (\tfrac12, 0)$ and $E(\pi) = (-\tfrac12, 0)$, the two
cusps. Setting $y = 0$ in the line equation: $x\sin2\theta =
\sin\theta$, so $x = \frac1{2\cos\theta} \to \frac12$ as
$\theta \to 0$: paraxial rays focus at distance $R/2$ from
the center --- the focal \omterm{def:b2:curves:length}{length} of the spherical mirror.

\textbf{20.} $\norm{E'(\theta)} = \tfrac32\abs{\sin\theta}$,
so the \omterm{def:b2:curves:length}{length} is $\tfrac32\int_0^{2\pi}\abs{\sin\theta}\,
\dd\theta = \tfrac32\cdot4 = 6$. Near $\theta = 0$, with
$\cos k\theta = 1 - \tfrac{k^2\theta^2}2 + O(\theta^4)$ and
$\sin k\theta = k\theta - \tfrac{k^3\theta^3}6 +
O(\theta^5)$:
\[
x - \tfrac12 = \frac{3\theta^2 + O(\theta^4)}{4}
= \tfrac34\theta^2 + O(\theta^4),
\qquad
y = \frac{4\theta^3 + O(\theta^5)}{4}
= \theta^3 + O(\theta^5) :
\]
$p = 2$, $q = 3$, a cusp of the first kind pointing along the
axis --- the bright point of the coffee-cup caustic.

\textbf{21.} Parametrize the illuminated points by
$\Phi(\theta, r) = P_\theta + r\,v_\theta$ (position along
each reflected ray). The Jacobian \omterm{def:b2:linalg:det}{determinant}
$\det(\partial_\theta\Phi, \partial_r\Phi) =
\det(P_\theta' + rv_\theta', v_\theta)$ is \omterm{def:b2:affine:subspace}{affine} in $r$ and
vanishes for exactly one $r = r_*(\theta)$ --- and
$\Phi(\theta, r_*(\theta))$ is the characteristic point, since
there the ray direction and the variation of the family become
dependent. Off the \omterm{pb:b2:curves:1}{envelope} the map is a local diffeomorphism,
and a point just inside the caustic is hit by two nearby rays
(two solutions $\theta$), a point outside by none from that
part of the family; on the caustic the two merge. Light
intensity is inversely proportional to the Jacobian's absolute
value, so it blows up along the \omterm{pb:b2:curves:1}{envelope}: the caustic is the
bright curve, brightest of all at the cusp, where the
degeneracy is worst.

\textbf{22.} $x' = 1 - \cos t$, $y' = \sin t$, $x'' = \sin
t$, $y'' = \cos t$, so $x'y'' - y'x'' = \cos t - 1$ and
$\norm{\gamma'}^2 = 2(1 - \cos t) = 4\sin^2\tfrac t2$;
by \cref{prop:b2:curves:kappaformula},
\[
\kappa(t) = \frac{-(1 - \cos t)}{8\sin^3\tfrac t2}
= \frac{-2\sin^2\tfrac t2}{8\sin^3\tfrac t2}
= -\frac1{4\sin\tfrac t2}
\qquad (0 < t < 2\pi).
\]
With $T = (\sin\tfrac t2, \cos\tfrac t2)$ (divide $\gamma'$
by $2\sin\tfrac t2$) and $N = (-\cos\tfrac t2, \sin\tfrac
t2)$:
\[
\gamma + \frac1\kappa N
= \gamma - 4\sin\tfrac t2\,\Bigl(-\cos\tfrac t2,\
\sin\tfrac t2\Bigr)
= \bigl(t - \sin t + 2\sin t,\ 1 - \cos t - 2(1 - \cos
t)\bigr),
\]
i.e.\ $c(t) = (t + \sin t,\ \cos t - 1)$. Substituting $t = u
+ \pi$:
\[
c = \bigl(u + \pi - \sin u,\ -\cos u - 1\bigr)
= \bigl((u - \sin u) + \pi,\ (1 - \cos u) - 2\bigr) :
\]
the cycloid $\gamma(u)$ translated by $(\pi, -2)$. The
evolute of a cycloid is a congruent cycloid, hanging one
level below.

\textbf{23.} $R(t) = 1/\abs{\kappa(t)} = 4\sin\tfrac t2$, so
$R(\pi) = 4$: half of the arch \omterm{def:b2:curves:length}{length} $8$ computed in
\cref{exo:b2:curves:1}. The evolute's velocity $c'(t) = (1 +
\cos t, -\sin t)$ vanishes at $t = \pi$: the cusp is $c(\pi)
= (\pi, -2)$, directly below the apex $\gamma(\pi) = (\pi,
2)$, at distance $4 = R(\pi)$, exactly the \omterm{def:b2:curves:length}{length} of the
osculating radius there.

\textbf{24.} From \cref{exo:b2:curves:6}, $c'(s) =
-\frac{\kappa'(s)}{\kappa(s)^2}N(s) = R'(s)\,N(s)$, so
$\norm{c'(s)} = \abs{R'(s)}$ and, for monotone $R$,
\[
\int_{s_0}^{s_1}\norm{c'(s)}\,\dd s =
\Bigl|\int_{s_0}^{s_1}R'(s)\,\dd s\Bigr| = \abs{R(s_1) -
R(s_0)} .
\]
Say $R$ decreases. A string laid along the evolute beyond
$c(s_0)$ and prolonged by the segment from $c(s_0)$ to
$\gamma(s_0)$ (which is tangent to the evolute, by question 4)
has, when peeled off up to $c(s)$ and stretched taut, straight
part of \omterm{def:b2:curves:length}{length} $R(s_0) - \bigl(R(s_0) - R(s)\bigr) = R(s)$
pointing from $c(s)$ along the normal --- landing exactly on
$\gamma(s)$: the free end traces the original curve
(``involute''). Huygens hung a pendulum between two cycloidal
cheeks: the cord wraps on the evolute, so the bob describes a
cycloid --- the tautochrone, whose oscillation period does
not depend on the amplitude.

\textbf{25.} Tangents of the parabola: $\Delta = 1$ and $E' =
(1, t) \neq 0$ everywhere --- smooth \omterm{pb:b2:curves:1}{envelope} (the parabola
itself). Sliding ladder: $\Delta = -1$, but $E' =
\tfrac32\sin 2t\,(-\cos t, \sin t)$ vanishes at the quadrant
ends --- the four cusps of the \omterm{pb:b2:curves:1}{astroid}. Normals of the
parabola and of the cycloid: $\Delta = \kappa \neq 0$, and
$E' = R'N$ vanishes exactly where the \omterm{thm:b2:curves:frenet2d}{curvature} is extremal
--- cusps of the evolutes at $(0,1)$ and $(\pi, -2)$. Caustic:
$\Delta = 2$, and $E' = \tfrac32\sin\theta\,(\cos2\theta,
\sin2\theta)$ vanishes at $\theta = 0, \pi$ --- the two cusps
of the \omterm{pb:b2:curves:1}{nephroid}, the focal points of the mirror. Extrema of
\omterm{thm:b2:curves:frenet2d}{curvature} \emph{must} produce cusps on the \omterm{pb:b2:curves:1}{envelope} of the
normals, since the evolute's velocity is $R'N$: this is why
the evolute of the ellipse has four cusps (four vertices),
and the degenerate families (pencil, parallels) are the cases
where the machine outputs a point or nothing at all.
\end{solution}
